\documentclass[14pt,a4paper]{extreport}
% preamble.tex — спільна преамбула для всіх лабораторних робіт. Підключається одразу після \documentclass у кожному lab*/main.tex.
%
% ЗБІРКА: LuaLaTeX (не pdfLaTeX). Причина: пакет listings під pdfLaTeX не читає UTF-8 і зупиняє збірку на кириличних коментарях у лістингах програмного коду. LuaLaTeX обробляє UTF-8 нативно.
%
% ШРИФТ: Liberation Serif — метрично сумісний з Times New Roman і містить повний набір кириличних гліфів (TeX Gyre Termes кирилиці не має, через що український текст не відображався взагалі).

\usepackage{fontspec}
\setmainfont{Liberation Serif}
\setmonofont{DejaVu Sans Mono}[Scale=0.82]

\usepackage{polyglossia}
\setmainlanguage{ukrainian}
\setotherlanguage{english}

\usepackage{microtype}

% Поля, інтервал, відступи й заголовки нижче взяті з шаблону НУЛП, за яким написана бакалаврська робота (../bachelor-thesis/thesis/main.tex). Той самий формат приймає В.М. Шиманський і для магістерських, тож усі звіти тримаємо в ньому, а не в тому, що переказували на консультації.
\usepackage[left=25mm,right=15mm,top=15mm,bottom=23mm,footskip=1.5mm]{geometry}
\usepackage{setspace}
% Саме \setstretch{1.5}, а не \onehalfspacing: останній дає близько 1.25.
\setstretch{1.5}

\usepackage{graphicx}
\usepackage{float}
\usepackage{booktabs}
\usepackage{array}
\usepackage{tabularx}
\usepackage{longtable}
\usepackage{enumitem}
% labelindent вирівнює маркер списку з абзацним відступом — так у шаблоні НУЛП.
\setlist{nosep}
\setlist[itemize]{label=\textbullet,leftmargin=*,labelindent=1.27cm}
\setlist[enumerate]{leftmargin=*,labelindent=1.27cm}
\usepackage{indentfirst}
\usepackage{amsmath}

\usepackage{hyperref}
\hypersetup{colorlinks=true,linkcolor=black,urlcolor=blue,citecolor=black}

% xurl дозволяє \url розривати адресу в будь-якому місці, а не лише після скісної риски. Без нього довгий сегмент шляху (напр. «scopus---podrobnaya- instruktsiya-dlya-uchenogo») виходить за межі сторінки, бо дефіс для hyperref не є точкою розриву. Вантажиться після hyperref — перевизначає його \UrlBreaks.
\usepackage{xurl}

% ---------------------------- Лістинги програмного коду ---------------------------- Використано minted (через Pygments), а не listings: listings не вміє коректно обробляти багатобайтові символи й «висмикує» кириличні коментарі з коду, зліплюючи їх в один рядок. Потребує прапорця -shell-escape.
\usepackage{xcolor}
\usepackage[newfloat]{minted}
% labelsep=period дає «Рис. 2.1. Назва» замість «Рис. 2: Назва» — формат НУЛП.
\usepackage[labelsep=period]{caption}

% baselinestretch=1 обов'язковий: без нього лістинг успадковує півторашку основного тексту, і довгий фрагмент коду перестає вміщатися у сторінку. Усередині [H]-флоата це не помилка збірки, а нескінченний цикл порожніх сторінок — lab1 так намолотив 161 тисячу сторінок.
\setminted{
  baselinestretch=1,
  fontsize=\footnotesize,
  linenos=true,
  frame=single,
  framesep=4pt,
  breaklines=true,
  breakanywhere=true,
  tabsize=4,
  numbersep=7pt,
  bgcolor=gray!8,
}
\setminted[python]{python3=true}

\SetupFloatingEnvironment{listing}{name=Лістинг,listname=Перелік лістингів}

% ---------------------------- Нумерація та оформлення заголовків ----------------------------
\usepackage{titlesec}

\setcounter{secnumdepth}{3}
\renewcommand{\thesection}{\arabic{section}}
\renewcommand{\thesubsection}{\arabic{section}.\arabic{subsection}}
\renewcommand{\thesubsubsection}{\arabic{section}.\arabic{subsection}.\arabic{subsubsection}}

% Розділ по центру, 16 пт, великими літерами; слово «РОЗДІЛ» опущене, бо звіт не дисертація. Номер лишається: \numberwithin прив'язує до нього нумерацію рисунків («Рис. 2.1»), і без нумерованих розділів вона б зламалася. \MakeUppercase четвертим аргументом отримує текст заголовка, тому в джерелі заголовки пишуться звичайним регістром.
\titleformat{\section}[block]
  {\centering\bfseries\fontsize{16}{19.2}\selectfont}
  {\thesection.}{0.3em}{\MakeUppercase}
% Підрозділи — зліва з абзацного відступу, звичайним регістром.
\titleformat{\subsection}[block]{\bfseries}{\thesubsection.}{1em}{}
\titleformat{\subsubsection}[block]{\bfseries\itshape}{\thesubsubsection.}{1em}{}

% Незіркові \titlespacing — обов'язкові. Без них titlesec застосовує зірковий варіант, який прибирає абзацний відступ у першого абзацу після заголовка й перекриває indentfirst. Значення — типові для report.
\titlespacing{\section}{0pt}{20pt}{6pt}
\titlespacing{\subsection}{1.27cm}{6pt}{0pt}
\titlespacing{\subsubsection}{1.27cm}{6pt}{0pt}

\setlength{\parindent}{1.27cm}
\setlength{\parskip}{0pt}
\renewcommand{\arraystretch}{1.3}

% Переноси слів вимкнені — так у шаблоні НУЛП. Натомість абзац розганяється міжслівними пробілами, тому tolerance і emergencystretch підняті: без них рядки виходили б за праве поле.
\hyphenpenalty=10000
\exhyphenpenalty=10000
\tolerance=9999
\setlength{\emergencystretch}{3em}

% ---------------------------- Підписи до рисунків і таблиць ----------------------------
\captionsetup[figure]{name=Рис.,justification=centering,font={bf,it}}
\captionsetup[table]{name=Таблиця,justification=raggedleft,singlelinecheck=false,labelfont=bf,textfont=normalfont}

% Рисунки, таблиці й формули нумеруються в межах розділу: «Рис. 2.1», а не «Рис. 5». Так у шаблоні НУЛП, і так само виглядали приклади на консультації.
\numberwithin{equation}{section}
\numberwithin{figure}{section}
\numberwithin{table}{section}

% ---------------------------- Колонтитули ----------------------------
\usepackage{fancyhdr}
\pagestyle{fancy}
\fancyhf{}
\fancyfoot[R]{\thepage}
\renewcommand{\headrulewidth}{0pt}
\fancypagestyle{plain}{%
  \fancyhf{}
  \fancyfoot[R]{\thepage}
  \renewcommand{\headrulewidth}{0pt}
}

% titlepage.tex — титульний аркуш, спільний для всіх дисциплін.
% Потребує визначених у subject.tex команд:
%   \SubjectName, \LecturerName, \StudentName, \StudentGroup, \AcademicYear
%
% Використання у тексті: \MakeLabTitle{1}{Тема роботи}

\newcommand{\MakeLabTitle}[2]{%
\begin{titlepage}
\begin{center}
  {\large\bfseries МІНІСТЕРСТВО ОСВІТИ І НАУКИ УКРАЇНИ}\\[6pt]
  {\large\bfseries Національний університет «Львівська політехніка»}\\[6pt]
  {\large\bfseries Інститут комп'ютерних наук та інформаційних технологій}\\[6pt]
  {\large\bfseries Кафедра «Системи штучного інтелекту»}\\[1.2cm]

  \includegraphics[width=5cm]{../../../../logo}\\[1.2cm]

  {\Large\bfseries ЛАБОРАТОРНА РОБОТА №#1}\\[0.5cm]
  з дисципліни\\[4pt]
  «\SubjectName»\\[0.6cm]
  на тему\\[4pt]
  \textit{#2}\\[2cm]
\end{center}

\begin{flushright}
  {\bfseries Виконав:}\\
  студент групи \StudentGroup\\
  \StudentName\\[1cm]
  {\bfseries Викладач:}\\
  \LecturerName\\
\end{flushright}

\vfill
\begin{center}
  Львів --- \AcademicYear
\end{center}
\end{titlepage}
\clearpage
}


% Метадані дисципліни. Використовуються титульним аркушем (titlepage.tex).

\newcommand{\SubjectName}{Штучний інтелект та інноваційні інформаційні технології}
% Два рядки: посада з науковим ступенем, далі прізвище. `\\` розриває рядок усередині flushright-блоку титульного аркуша.
\newcommand{\LecturerName}{доц. каф. СШІ к.т.н.,\\ Думин І.Б.}
\newcommand{\StudentName}{Галич Максим}
\newcommand{\StudentGroup}{КНСШ-13}
\newcommand{\AcademicYear}{2026}

% Рисунки — згенеровані артефакти, тому лежать не поряд із текстом,
% а в .cache/ (див. ARTIFACTS у src/config.py). Імена всередині
% (eda/..., lab1/...) не змінилися.
\graphicspath{{../../../../.cache/reports-artifacts/1th term 2026 Autumn/Штучний інтелект та інноваційні інформаційні технології/}}

\begin{document}

\MakeLabTitle{2}{Побудова та аналіз моделей штучного інтелекту з використанням PyTorch}

\textbf{Мета роботи.} Опанувати практику створення й оптимізації моделей
машинного навчання у PyTorch, навчитися застосовувати різні архітектури,
регуляризацію та методи покращення узагальнюючої здатності моделі, виконати
порівняльний аналіз результатів, зокрема порівняти PyTorch із TensorFlow.

\section{Постановка експерименту}

Лабораторна робота №2 відтворює засобами PyTorch той самий експеримент, що
був виконаний у лабораторній роботі №1 засобами TensorFlow. Використано той
самий набір даних \textbf{California Housing} (20\,640 записів, 8 ознак,
задача регресії), той самий розвідувальний аналіз і ту саму попередню обробку.

Щоб порівняння фреймворків було коректним, усі параметри експерименту винесено
у спільний модуль \texttt{src/config.py}, який імпортують обидві реалізації.
Таким чином, будь-яка різниця в результатах зумовлена фреймворком, а не
розбіжністю в налаштуваннях.

% Без оточення listing: цей лістинг (увесь config.py) вищий за сторінку, а
% плаваючий об'єкт розривати не можна — [H] не має куди його покласти і збірка
% зациклюється, друкуючи порожні сторінки без жодної помилки. Звичайний
% \inputminted переноситься між сторінками, \captionof дає ту саму нумерацію.
\inputminted{python}{../src/aiit/config.py}
\captionof{listing}{Спільні гіперпараметри обох лабораторних робіт (\texttt{src/config.py})}
\label{lst:config}

\begin{table}[H]
\caption{Умови експерименту, однакові для обох фреймворків}
\label{tab:setup}
\begin{tabularx}{\linewidth}{|X|l|}
\hline
\textbf{Параметр} & \textbf{Значення} \\
\hline
Поділ вибірки & 60/20/20 зі стратифікацією \\
\hline
Оптимізатор & Adam \\
\hline
Швидкість навчання & $10^{-3}$ \\
\hline
Розмір батча & 256 \\
\hline
Кількість епох & 20 \\
\hline
Функція втрат & MSE \\
\hline
Зерно генератора & 42 \\
\hline
Обчислювальний пристрій & CPU \\
\hline
\end{tabularx}
\end{table}

\section{Реалізація архітектур у PyTorch}

Реалізовано ті самі дві архітектури: базову feedforward мережу з двома
прихованими шарами та ускладнену мережу з Dropout, Batch Normalization
і L2-регуляризацією. Кожну архітектуру навчено з активаціями ReLU та GELU.

Принципова відмінність від Keras полягає у способі задання L2-регуляризації.
У Keras вона вказується на рівні шару через \texttt{kernel\_regularizer},
натомість у PyTorch її задає оптимізатор параметром \texttt{weight\_decay}.

\begin{listing}[H]
\inputminted[firstline=26,lastline=52]{python}{../src/aiit/torch_models.py}
\caption{Архітектури моделей у PyTorch (\texttt{src/torch\_models.py})}
\label{lst:ptmodels}
\end{listing}

\section{Цикл навчання}

Друга принципова відмінність --- відсутність у PyTorch готового методу
\texttt{.fit()}. Цикл навчання доводиться писати вручну: перебір епох і
батчів, обнулення градієнтів, прямий прохід, обчислення втрати, зворотне
поширення та крок оптимізатора.

Окремої уваги потребує перемикання режимів \texttt{model.train()} та
\texttt{model.eval()}. Без нього Dropout лишався б активним під час
обчислення валідаційної втрати, а BatchNorm використовував би статистики
поточного батча замість накопичених, що спотворило б результати.

\begin{listing}[H]
\inputminted[firstline=57,lastline=89]{python}{../src/aiit/run_lab2.py}
\caption{Ручний цикл навчання у PyTorch (\texttt{src/run\_lab2.py})}
\label{lst:ptloop}
\end{listing}

\section{Результати PyTorch}

\subsection{Метрики якості}

\begin{table}[H]
\caption{Метрики моделей на тестовій вибірці (PyTorch)}
\label{tab:ptresults}
\begin{tabularx}{\linewidth}{|X|r|r|r|r|r|}
\hline
\textbf{Модель} & \textbf{RMSE} & \textbf{MAE} & \textbf{$R^2$} &
\textbf{MAPE, \%} & \textbf{Час, с} \\
\hline
baseline-relu       & \textbf{0,6165} & \textbf{0,4313} & \textbf{0,7114} & 24,0 & 4,20 \\
\hline
baseline-gelu       & 0,6228 & 0,4322 & 0,7055 & 23,8 & 6,12 \\
\hline
regularized-gelu    & 0,7495 & 0,4948 & 0,5734 & 23,2 & 15,72 \\
\hline
regularized-relu    & 0,7542 & 0,4952 & 0,5681 & 22,9 & 4,56 \\
\hline
\end{tabularx}
\end{table}

Як і у TensorFlow, найкращою виявилася базова мережа з активацією ReLU:
RMSE $= 0{,}617$, $R^2 = 0{,}711$.

\subsection{Навчальні криві}

\begin{figure}[H]
  \centering
  \includegraphics[width=\linewidth]{lab2/learning_curves}
  \caption{PyTorch: криві тренувальної та валідаційної втрати}
  \label{fig:ptcurves}
\end{figure}

\begin{figure}[H]
  \centering
  \includegraphics[width=0.8\linewidth]{lab2/overfit_gap}
  \caption{PyTorch: розрив між валідаційною та тренувальною втратою}
  \label{fig:ptgap}
\end{figure}

\begin{figure}[H]
  \centering
  \includegraphics[width=0.62\linewidth]{lab2/pred_vs_actual}
  \caption{PyTorch: прогнозовані значення проти фактичних (baseline-relu)}
  \label{fig:ptpred}
\end{figure}

\section{Порівняння TensorFlow та PyTorch}

\subsection{Якість моделей}

\begin{table}[H]
\caption{Порівняння RMSE на тестовій вибірці}
\label{tab:cmp}
\begin{tabularx}{\linewidth}{|X|r|r|r|}
\hline
\textbf{Модель} & \textbf{TensorFlow} & \textbf{PyTorch} & \textbf{Різниця} \\
\hline
baseline-relu     & 0,6082 & 0,6165 & $+1{,}4$~\% \\
\hline
baseline-gelu     & 0,6181 & 0,6228 & $+0{,}8$~\% \\
\hline
regularized-relu  & 0,7480 & 0,7542 & $+0{,}8$~\% \\
\hline
regularized-gelu  & 0,7770 & 0,7495 & $-3{,}5$~\% \\
\hline
\end{tabularx}
\end{table}

\begin{figure}[H]
  \centering
  \includegraphics[width=\linewidth]{lab2/framework_comparison}
  \caption{Порівняння якості та швидкодії TensorFlow і PyTorch}
  \label{fig:cmp}
\end{figure}

Розбіжності між фреймворками не перевищують 3,5~\% і не мають систематичного
характеру: у трьох випадках дещо кращим виявився TensorFlow, в одному ---
PyTorch. Це узгоджується з очікуваннями: за ідентичних архітектури,
гіперпараметрів і поділу даних обидва фреймворки реалізують один і той самий
алгоритм, а залишкова різниця пояснюється різною ініціалізацією ваг та різним
порядком формування батчів, які неможливо зробити побітово однаковими через
різні генератори випадкових чисел.

\subsection{Швидкодія}

На CPU час навчання 20 епох склав 2,3--8,0 с для TensorFlow та 4,2--15,7 с
для PyTorch. Виміри помітно зашумлені через конкуренцію за ядра процесора,
тому робити з них висновок про принципову перевагу котрогось фреймворку
некоректно. Можна лише констатувати, що на задачі такого масштабу
(12 тис. прикладів, 2,7 тис. параметрів) обидва фреймворки дають цілком
прийнятний час навчання без використання GPU.

\subsection{Відмінність у підрахунку параметрів}

Окремо слід пояснити розбіжність у кількості параметрів регуляризованих
моделей: 2\,977 у TensorFlow проти 2\,785 у PyTorch. Архітектури ідентичні,
різниця складає рівно $2 \times (64 + 32) = 192$ значення.

Причина суто облікова. Шар BatchNorm має чотири набори величин: навчувані
$\gamma$ і $\beta$ та ненавчувані ковзні середнє й дисперсію. Метод
\texttt{model.count\_params()} у Keras враховує всі чотири, тоді як
\texttt{sum(p.numel() for p in model.parameters())} у PyTorch повертає лише
навчувані, бо ковзні статистики зареєстровані як буфери, а не параметри.
Отже, \textbf{моделі однакові}, різниться лише спосіб підрахунку.

\subsection{Переваги та недоліки фреймворків}

\begin{table}[H]
\caption{Порівняння фреймворків за досвідом виконання роботи}
\label{tab:prosandcons}
\begin{tabularx}{\linewidth}{|l|X|X|}
\hline
\textbf{Критерій} & \textbf{TensorFlow / Keras} & \textbf{PyTorch} \\
\hline
Обсяг коду &
Менший: \texttt{.fit()} інкапсулює весь цикл навчання &
Більший: цикл навчання пишеться вручну \\
\hline
Контроль над навчанням &
Нестандартна логіка потребує callback-ів або \texttt{train\_step} &
Повний контроль за замовчуванням \\
\hline
Регуляризація L2 &
На рівні шару (\texttt{kernel\_regularizer}) &
На рівні оптимізатора (\texttt{weight\_decay}) \\
\hline
Режими train/eval &
Перемикаються автоматично &
Вимагають явних \texttt{.train()} / \texttt{.eval()} \\
\hline
Імовірні помилки &
Приховування деталей ускладнює діагностику &
Легко забути \texttt{zero\_grad()} або \texttt{.eval()} \\
\hline
Налагодження &
Складніше через рівень абстракції &
Простіше: звичайний імперативний код Python \\
\hline
\end{tabularx}
\end{table}

\section{Аналіз архітектур та причин недонавчання}

Обидва фреймворки дали узгоджений і спершу неочікуваний результат:
\textbf{базова мережа переважає регуляризовану} приблизно на 22~\% за RMSE.

Діагностика за рисунками~\ref{fig:ptcurves} та~\ref{fig:ptgap} показує, що
причиною є \textbf{недонавчання, а не перенавчання}. Розрив між валідаційною
та тренувальною втратою наприкінці навчання не перевищує 0,081 і не зростає
з часом, тобто перенавчання відсутнє. Натомість регуляризовані моделі
демонструють високу
втрату \textit{одночасно} на тренувальній та валідаційній вибірках
(0,472 і 0,553 проти 0,342 і 0,362 у базової моделі) --- класична ознака
надмірного обмеження моделі.

Пояснення полягає у співвідношенні ємності моделі та обсягу даних. Мережа
має близько 2,7 тис. параметрів на 12,4 тис. тренувальних прикладів, тобто
даних майже вп'ятеро більше, ніж параметрів. За такого співвідношення
перенавчання малоймовірне, і регуляризувати нічого. Dropout із коефіцієнтом
0,2 вимикає кожен п'ятий нейрон у шарах із 64 та 32 одиницями, що для мережі
такого розміру є відчутною втратою ємності.

\textbf{Висновок для практики:} регуляризацію слід вводити як відповідь на
діагностоване перенавчання, а не як обов'язковий елемент архітектури.
Застосована «про всяк випадок», вона погіршує результат.

\section{Висновки}

У ході лабораторної роботи засобами PyTorch відтворено експеримент
лабораторної роботи №1 та виконано порівняльний аналіз двох фреймворків.

\begin{enumerate}
\item Реалізовано дві архітектури (базову та регуляризовану) з двома
функціями активації кожна. Найкращий результат показала базова мережа
з ReLU: RMSE $= 0{,}617$, $R^2 = 0{,}711$.

\item Результати PyTorch і TensorFlow збігаються в межах 3,5~\% і не
виявляють систематичної переваги жодного з фреймворків. Залишкова
розбіжність пояснюється різною ініціалізацією ваг та порядком формування
батчів.

\item Виявлену розбіжність у кількості параметрів (2\,977 проти 2\,785)
пояснено способом підрахунку: Keras враховує ковзні статистики BatchNorm,
PyTorch --- ні. Самі моделі ідентичні.

\item Основна практична відмінність фреймворків --- рівень абстракції.
Keras приховує цикл навчання за \texttt{.fit()}, що скорочує код, але
ускладнює нестандартні сценарії. PyTorch вимагає написати цикл вручну,
що збільшує обсяг коду та кількість потенційних помилок, але дає повний
контроль і простіше налагодження.

\item В обох фреймворках регуляризація погіршила результат приблизно на
22~\% за RMSE. Встановлено, що причина --- недонавчання: обсяг даних
переважає ємність моделі майже вп'ятеро, тому перенавчання не виникає
і регуляризація лише обмежує апроксимаційну здатність мережі.

\item Для цієї задачі вибір архітектури впливає на результат значно сильніше
(різниця 22~\%), ніж вибір функції активації (до 1,0~\% у PyTorch та до
3,9~\% у TensorFlow) чи вибір фреймворку (до 3,5~\%).
\end{enumerate}

\end{document}
